\section{Preliminaries }
\label{sec:preliminaries}
In the present section we show that the construction of type B Fock space \cite{BozejkoEjsmontHasebe2015} can be adapted to recover the
double Fock space of type B.
In the following we will briefly describe the tools we use in this investigation,
which is partially contained in the previous paper \cite{BozejkoEjsmontHasebe2015}.
For further information the reader is referred to \cite{BozejkoEjsmontHasebe2015} and the references therein.

\subsection{Coxeter groups of type B}
Recall that the Coxeter
group of type B (also known as the hyperoctahedral group) of degree $n$, denoted by $B(n)$, is generated
by the elements $\pi_0, \pi_1, \dots , \pi_{n-1}$ subject to the defining relations $\pi_i^2=e$, $0\leq i \leq n-1$, $(\pi_0\pi_1)^4=(\pi_i \pi_{i+1})^3=e$, $1\leq i < n-1$ and $(\pi_i \pi_j)^2=e$ if $|i-j|\geq2$, $0\leq i,j\leq n-1$. Note that $\{\pi_i\mid i=1,\dots,n-1\}$ generate the symmetric group $S(n)$.
The Coxeter diagram for $B(n)$ is described in Figure \ref{fig:BN}.
\begin{figure}[htp]
\centering
		\begin{tikzpicture}
			[scale=.5,auto=left,every node/.style={circle}]
			\node (n7) at (1,3.6) {$\pi_{n-1}$};
			\node (n6) at (1,3) {$\circ$};
			\node (n5) at (-3,3) {$\dots$};
			\node (n3) at (-7,3.6) {$\pi_2$};
			\node (n1) at (-7,3) {$\circ$};
			\node (n4) at (-11,3.6) {$\pi_{1}$};
			\node (n2) at (-11,3) {$\circ$};
			\node (n8) at (-15,3.6) {$\pi_{0}$};
			\node (n0) at (-15,3) {$\circ$};
			
			\foreach \from/\to in {n2/n1, n1/n2,n1/n5,n5/n6}
			\draw (\from) -- (\to);
			\draw (-14.5,3.1) -- (-11.5,3.1);
			\draw (-14.5,2.9) -- (-11.5,2.9);
		\end{tikzpicture}
		\caption{The Coxeter diagram for $B(n)$.}
		\label{fig:BN}
\end{figure}


We express $\sigma \in B(n)$ in an irreducible form
\[
\sigma=\pi_{i_1} \cdots \pi_{i_{k}}, \qquad 0\leq i_1,\dots,i_k \leq n-1,
\]
i.e., in a form with the minimal length, and in this case let
\begin{align*}
	&l_1(\sigma)= \text{the number of the occurrences of the factor $\pi_0$ in $\sigma$},
\\[1mm]
	&l_2(\sigma) = \text{the number of the occurrences of all the factor of the form $\pi_1,\dots,\pi_{n-1}$ in $\sigma$}.
\end{align*}

\begin{Remark}
These definitions do not depend on the way we express $\sigma$ in an irreducible form, and therefore, $l_1(\sigma)$ and $l_2(\sigma)$ are well defined (see \cite[Proposition~1]{BozejkoSzwarc2003}).
\end{Remark}
